\subsection{Coordinates defined by the actual passages}
\label{rw:coordinates}

Coefficient paths must control both the slow factors and their moving
entries on the original sections. We define the gaps and splittings
from the actual weak graphs and strong fibres, then obtain common
inverse boxes at fixed normalization. To retain order $N$ through the
minimum, Morse coordinate and time unit, reserve $N+4$ reduced-phase
orders before choosing the larger domains. Connections to $Y=0$ use
Euler scale derivatives.

We use the following elementary uniform inverse construction repeatedly.
\begin{lemma}[Common inverse boxes]\label{rw:coordinates:inverse-box}
Let $F_\epsilon$ be a family of coefficient maps on a fixed real
neighborhood, with bounded derivatives through a prescribed order,
uniformly for sufficiently small $\epsilon\ge0$, and suppose
$F_\epsilon\to F_0$ in the coefficient $C^1$ norm. Suppose its Jacobian
at a base point and $\epsilon=0$ is invertible. After reducing the
neighborhood, there are nested fixed product boxes $Q\Subset Q^+$ in
the new coordinates on which all $F_\epsilon$ have single-valued
inverses with bounded finite jets. Any additional passive coefficients
can be included in these boxes. The assertion remains valid with Euler
scale jets when that is the regularity of $F_\epsilon$.
\end{lemma}
\begin{proof}
Write $A=DF_0$ at the base point. On a small convex source ball arrange
$\|A^{-1}(DF_\epsilon-A)\|<1/2$, uniformly in the passive parameters
and scale. The equation $F_\epsilon(x)=y$ is then the contraction
\[
 x=x_*+A^{-1}\{y-F_\epsilon(x_*)
              -[F_\epsilon(x)-F_\epsilon(x_*)-A(x-x_*)]\}.
\]
Choose a common target ball whose contractions stay in the source
ball, and nested product boxes $Q,Q^+$ inside it. The inverse
Jacobian is bounded by $2\|A^{-1}\|$; differentiating
$F_\epsilon\circ F_\epsilon^{-1}=\mathrm{id}$ gives all higher
jets inductively, including passive and Euler derivatives.
\end{proof}

\begin{proposition}[Actual gap and splitting coordinates]
\label{rw:center-coordinates}
On the joint repeated and weak blocks specified below, the actual
passages admit the following coefficient coordinates on common nested
boxes, with uniformly bounded selected finite inverse jets:
\[
 (\bar B,\delta,\tau,\zeta,\nu),\qquad
 (\bar B,\bar c,d_s,\varsigma,\nu),
\]
respectively. The joint gaps are exactly $\delta-\zeta$ and
$\delta+\zeta$; their center axis is $\tau=\zeta=\nu=0$.
The weak splitting is $d_s=P_L-\widehat p$, and its center axis is
$d_s=\varsigma=\nu=0$. The joint entries can be calibrated on the
original finite section so that both equal $\Xi-\bar Bp$ on the
center and have bounded coefficient Lipschitz constants. On the weak
block the regular horizontal input $\chi(u)$ has bounded coefficient
jets, while the exact moving eigenvalue $\gamma$ and the composed
horizontal factor $F_I(\chi(u))$ have the positive center $\bar B$
derivatives stated below. All these constructions precede
comparison of coefficient paths or zeros of a displacement.
\end{proposition}
\begin{proof}
\leavevmode\par\smallskip\noindent\textit{The joint repeated block.}
Fix $g_*>0$ and the limiting coefficients
\[
 \bar B=g_*^2,\qquad \bar m=-g_*^2,\qquad
 \bar a=-2g_*,\qquad \bar c=2g_*,\qquad \nu=\epsilon=0.
\]
Write $b=1-\epsilon\bar B$ and keep $\bar B$ and $b$ bounded away
from zero. The upper gap is the exact quantity already prepared in
\eqref{rw:upper:repeated-normal}:
\begin{equation}\label{rw:coordinates:upper-gap}
 \mu_U=-\frac{(\bar c+\nu)^2+4b\bar m}{4b^2}.
\end{equation}
For the lower passage we select the phase within the freedom of the
construction in Proposition~\ref{rw:phase}. The following calibration
both specifies its gap and identifies the center input.

On the second center $\bar a=-\bar c$, $\bar m=-\bar B$, $\nu=0$,
define
\begin{equation}\label{rw:coordinates:center-quotient}
 \Xi=\frac{\bar c(1-2\epsilon\bar B)}{2b},\qquad
 \delta_c=\frac{4\bar Bb-\bar c^2}{4b^2},\qquad
 V=\frac{\Xi/\bar B-X-\bar cY/(2b)}{Y+1/\bar B}.
\end{equation}
Its denominator has a positive margin on each strip under
consideration. In the notation of the real center construction,
$V=\bar B\chi+\Xi$; alternatively direct differentiation in
\eqref{rw:geometry:upper-field} gives
\begin{equation}\label{rw:coordinates:center-scalar}
 V_\xi=bY(V^2+\delta_c),\qquad
 V(X,-\epsilon)=\frac{\bar c/2-\bar BX}{b},\qquad
 V(p,0)=\Xi-\bar Bp.
\end{equation}
These identities hold for positive, zero and negative $\delta_c$.

In the lower strip choose a fixed $0<z_h<1$ and reverse time.
Division by the positive factor $b(1-z)$ gives
$\dot V=\epsilon(V^2+\delta_c)$ on the center.
The normal coordinate $z$ contracts there: its linear rate on $z=0$
is $(X-\epsilon\bar c)/b$, with a positive margin near
$X=1/g_*$. In a smaller strip both this rate and the derivative
of $z$ times its rate remain positive. Extend the equations first
in these center coordinates, retaining the exact scalar equation and
normal contraction, and retaining $z=0$ invariance whenever $\nu=0$;
extend only the off-center defect by cutoffs.
This is an admissible choice in the graph and phase construction of
Section~\ref{rw:transitions}. It keeps the exact quotient on the
center and preserves the factor $\epsilon$ in the slow equation.
The latter also follows directly: $V_X=O(1)$, $V_z=O(\epsilon)$,
while the original $X$ velocity contains $\epsilon$.

The resulting actual lower weak equation is
$\dot w=\epsilon f_L^0(w;\mathbf b,\epsilon)$, with uniformly
bounded finite jets and $(f_L^0)_{ww}>1$. Its minimum $w_0$ is
therefore a finite-smooth coefficient function. Set
\begin{equation}\label{rw:coordinates:lower-gap}
 \mu_L=f_L^0(w_0),\qquad
 u=(w-w_0)\left[\int_0^1(1-s)
    (f_L^0)_{ww}(w_0+s(w-w_0))\,ds\right]^{1/2}.
\end{equation}
This is an increasing state coordinate, and $f_L^0=\mu_L+u^2$.
Dividing by its positive derivative gives exactly
$\dot u=\epsilon(\mu_L+u^2)$ while preserving the normal margin.
On the center, $w=V$, $w_0=0$ and $u=V$.

The normalization in \eqref{rw:coordinates:lower-gap} is fixed by the
preceding time and phase choice. To calculate it, take $\nu=0$.
The original line $z=0$ is invariant, and on it
$w=(\bar c/2-\bar BX)/b$. Restriction of the reversed field divided
by $b(1-z)$ gives, with $\tau=\bar a+\bar c$,
\begin{equation}\label{rw:coordinates:lower-exact}
 f_L^0(w)=w^2-\frac{\tau}{b}w
  +\frac{\bar B(1+\epsilon\bar m)+\bar a\bar c/2+\bar c^2/4}{b^2},
 \qquad
 \mu_L=\frac{\bar B(1+\epsilon\bar m)-\bar a^2/4}{b^2}.
\end{equation}
Thus $\mu_L$ is the finite-$X$ minimum multiplied by the positive
unit $\bar B/b^2$; the compact horizontal factor $F_I$ has a
different normalization. For $\nu\ne0$,
\eqref{rw:coordinates:lower-gap} continues to use the actual weak
graph and the chosen phase.

Set
\begin{equation}\label{rw:coordinates:joint}
 \delta=\frac{\mu_U+\mu_L}{2},\qquad
 \zeta=\frac{\mu_U-\mu_L}{2},\qquad
 \tau=\bar a+\bar c,\qquad \varsigma=\bar B+\bar m.
\end{equation}
At $\nu=0$ the exact identity
\begin{equation}\label{rw:coordinates:zeta-identity}
 \zeta=-\frac{\bar c\tau}{4b^2}
          +\frac{\tau^2}{8b^2}-\frac{\varsigma}{2b^2}
\end{equation}
follows by subtraction of the two gaps. At $\tau=\varsigma=0$,
\begin{equation}\label{rw:coordinates:joint-jacobian}
 \det\frac{\partial(\delta,\tau,\zeta)}
                 {\partial(\bar c,\tau,\varsigma)}
       =\frac{\bar c}{4b^4}>c_*>0.
\end{equation}
For example, on this axis $\delta_{\bar c}=-\bar c/(2b^2)$ and
$\zeta_{\varsigma}=-1/(2b^2)$, while
$\zeta_{\bar c}=0$, which proves the determinant. Continuity of the
actual phase and minimum derivatives retains a uniform nonzero margin
for both signs of small $\nu$ and small positive scale.

Consequently $(\bar B,\delta,\tau,\zeta,\nu)$ are coefficient
coordinates on common boxes as in
Lemma~\ref{rw:coordinates:inverse-box}. On these boxes
\begin{equation}\label{rw:coordinates:joint-center-axis}
 \tau=\zeta=\nu=0
 \quad\Longleftrightarrow\quad
 \bar a=-\bar c,\quad\bar m=-\bar B,\quad\nu=0.
\end{equation}
Indeed \eqref{rw:coordinates:zeta-identity} first gives
$\varsigma=0$ when $\tau=\zeta=\nu=0$; the converse follows from
the gap formulas. The center with prescribed $(\bar B,\delta)$ has
\begin{equation}\label{rw:coordinates:center-c}
 \bar c_c=2\sqrt{\bar Bb-b^2\delta}>0.
\end{equation}
All these inverse jets have uniform finite bounds on the larger box.
The two factors are exactly $u^2+\delta-\zeta$ and
$v^2+\delta+\zeta$. Thus fixed full passages across both minima
pass exactly when $\delta>|\zeta|$; moving entries use the whole-path tests
of Section~\ref{rw:transitions}.

\par\smallskip\noindent\textit{Entries on one initial section.}
The common boxes must also control the entry functions. At a finite
joint contact fix $p$ in a small interval about $P=1/g_*$ on $Y=0$.
Continue the outgoing lower branch backwards to $z=z_h$ and the upper
branch forwards to a fixed positive radial section. Both segments are
chosen before their slow passages. On the lower segment use $z$ as
the independent variable; $X$ has a positive margin and $z\ge z_h$,
so its equation and its finite parameter variations have bounded
coefficients. The upper segment first uses the low-height estimate
following \eqref{rw:upper:curvature}, then a compact regular arc.
Even at the repeated limiting input $p=P$ it reaches the fixed radial
section in finite affine time. The entry maps and their transverse
inverses have bounded ordinary $p$, coefficient and Euler scale jets
on a slightly larger common domain.

Choose the upper auxiliary phase to preserve
\eqref{rw:coordinates:center-quotient} on the center as well.
This agrees with its boundary coordinate at $R=0$ and with its exact
gap \eqref{rw:coordinates:upper-gap}. On the center the positive
upper time factor is $bY$, and the lower reversed time factor is
$b(1-z)$. Record their accumulated scalar times on the initial
segments, denoted $T_U$ and $T_L$. These have bounded selected jets:
the upper time integral has integrand $bY/(Y+\epsilon)$ times a
bounded regular factor, and the lower one is taken on a fixed
$z$-interval. The upper original transit time may contain
$|\log\epsilon|$, but this scalar time does not.

Let $u_e,v_e$ be the prepared entries at the fixed fast sections.
Apply the inverse scalar flows for these bounded recorded times to
define entries $e_L,e_U$ on the original initial section. Off the
center use the same smooth time integrals and the same complete scalar
extensions. Finite-time scalar flow and its inverse have bounded jets
on the larger box. Their roots and their through components are
invariant, so this change of entry does not alter the whole-path
passing tests in a smaller contact neighborhood. On the center, the
exact flow identities in \eqref{rw:coordinates:center-scalar} give
\begin{equation}\label{rw:coordinates:calibrated-entry}
 e_L(p)=e_U(p)=U_c(p):=\Xi-\bar Bp.
\end{equation}
The bounded regular segments therefore calibrate the entries before
the weak Dulac passage.

With $\mathbf q=(\bar B,\delta,\tau,\zeta,\nu)$ in the common
larger box, the chain rule and the inverse bounds just proved give,
for some fixed $L$,
\begin{equation}\label{rw:coordinates:entry-lipschitz}
 |e_i(p;\mathbf q_1)-e_i(p;\mathbf q_2)|
       \le L|\mathbf q_1-\mathbf q_2|,\qquad i=L,U.
\end{equation}
The bound also holds for each lower-order selected input jet when
one more coefficient derivative is reserved. Since
$\partial_pU_c=-\bar B$, both entry derivatives have a fixed
negative margin after reducing the transverse parameter box.
For $A=|\tau|+|\zeta|+|\nu|$, comparison at fixed
$(\bar B,\delta)$ yields
\begin{equation}\label{rw:coordinates:entry-center-error}
 |e_i(p;\mathbf q)-U_c(p;\bar B,\delta)|\le CA.
\end{equation}
The same construction on a separated finite-entry patch uses fixed
slow sections. At its radial end the common actual input of
Lemma~\ref{rw:transitions:common-input} supplies the transition jets.
The comparison paths below use these patchwise coordinate and entry
bounds to preserve an entire passing interval.

\par\smallskip\noindent\textit{The weak boundary with either sign of $\bar B$.}
Fix $\bar c_0>0$ and the limiting coefficients
$\bar B=\bar m=0$, $\bar a=-\bar c=-\bar c_0$.
Keep the exact $\gamma$ of \eqref{rw:upper:weak-parameters}; it is a
finite-smooth function on an ordinary coefficient neighborhood open in
both $\bar B$ and $\bar m$. Let $P_L$ be the actual outgoing lower
strong intersection with $Y=0$, and let $\widehat p$ be the actual
upper weak intersection used in \eqref{rw:upper:weak-rate}. Define
\begin{equation}\label{rw:coordinates:weak}
 d_s=P_L-\widehat p,\qquad \varsigma=\bar B+\bar m,
 \qquad\mathbf q_w=(\bar B,\bar c,d_s,\varsigma,\nu).
\end{equation}
These are actual section positions, before either one is changed to a
normalized power coordinate.

Choose the upper auxiliary graph to preserve the affine branch of
Lemma~\ref{rw:centers:weak-boundary}. In the present notation put
\begin{equation}\label{rw:coordinates:weak-center-branch}
 \Delta_b=\sqrt{\bar c^2-4b\bar B},\quad
 k_s=\frac{-\bar c-\Delta_b}{2b},\quad
 p_c=\frac{2(1-\epsilon\bar c^2)}
              {\bar c(1-2\epsilon\bar B)+\Delta_b}.
\end{equation}
Thus $p_c=\widehat p_C$ from
Lemma~\ref{rw:centers:weak-boundary}; its denominator has a positive
margin through $\bar B=0$. That lemma identifies the invariant line
$X=p_c+k_sY$ and both of its actual intersections with $Y=0$.
Subtract this known graph before making the auxiliary extension and
extend only its off-center defect. The graph transform then retains
this exact branch on the center for either sign of $\bar B$.

The lower simple-root preparation and the upper weak graph with its
finite connector supply the coefficient jets. At the zero-scale limiting
point the lower root solves $1+\bar aP_L+\bar BP_L^2=0$, while
the affine upper weak line has intersection $\widehat p=1/\bar c$
independent of $\bar a$. Implicit differentiation therefore gives
\begin{equation}\label{rw:coordinates:weak-jacobian}
 \partial_{\bar a}d_s\longrightarrow1/\bar c_0^2>0.
\end{equation}
The zero-scale limit holds in the coefficient jets on a fixed small
block: the only nonregular connector correction is the integrable
low-layer carrier $\epsilon/(Y+\epsilon)$. Hence this derivative
retains a fixed positive margin. Together with
$\partial_{\bar m}\varsigma=1$ and $\partial_{\bar a}\varsigma=0$,
this verifies the hypotheses of
Lemma~\ref{rw:coordinates:inverse-box} for the weak coefficient map.

The second center is exactly
\begin{equation}\label{rw:coordinates:weak-center-axis}
 d_s=\varsigma=\nu=0.
\end{equation}
One direction follows from the invariant branch. Conversely at
$\varsigma=\nu=0$ the actual center $\bar a=-\bar c$ gives
$d_s=0$; the nonzero splitting derivative and uniqueness in the
common inverse box show it is the only such $\bar a$.
On this center, direct substitution into the actual weak eigenvalue
in \eqref{rw:upper:weak-parameters} gives
\begin{equation}\label{rw:coordinates:gamma-center}
 \gamma_c=\bar Bp_c,\qquad
 \left.\partial_{\bar B}\gamma_c\right|_{\bar B=0}
   =\frac{1-\epsilon\bar c^2}{\bar c}>c_*>0.
\end{equation}
This is also the eigenvalue of the restricted equation
\eqref{rw:centers:weak-restricted}. Continuity retains its positive
derivative margin on a smaller common center block.

Finally we specify the regular connection used together with the
weak inverse map. Select the outgoing fast section of the reversed
weak block on the side $\varphi<\varphi_s$, a fixed small distance
from the weak point. Following the original forward field from this
section to the incoming horizontal section $\varphi=-1/v_h$ gives
a compact angular connection; choose $v_h$ small enough. At the
limiting field its angular velocity is
$-\varphi(\varphi+\bar c_0)<0$ on this closed arc, with a fixed
margin, and it retains that margin nearby. Both sections are
transverse. The radial equation is divisible by $R$, so its transfer
is $R$ times a positive unit with bounded jets and positive derivative
at zero. The weak input $u$ is likewise $R$ times a positive unit on
the first section. At the horizontal end $T=R/(-\varphi)$ and the
lower phase is a positive regular function of $T$. Their composition
therefore gives
\begin{equation}\label{rw:coordinates:weak-connector}
 x=\chi(u;\mathbf q_w,\epsilon),\qquad
 \chi(0)=0,\quad \chi_u>c_*>0,
\end{equation}
with bounded finite jets and inverse jets on a common signed
auxiliary input interval. It gives the horizontal entry before the
slow passage, so its regular inverse remains bounded even when
$t(u)$ is exponentially small.

Let $H(\mathbf q_w,u)=F_I(\chi(u;\mathbf q_w))$.
It and $\gamma(\mathbf q_w)$ have bounded coefficient derivatives
on the whole larger box. On the center, $\nu=0$ and
\eqref{rw:transitions:lower-factor} gives exactly
\[
 F_{I,c}(x)=\bar B+\bar cx+(1-\epsilon\bar B)x^2.
\]
Because $\chi_c(0)=0$ for all passive parameters, Hadamard integration
gives $\partial_{\bar B}\chi_c(u)=O(u)$, with the required input
jets. Consequently
\begin{equation}\label{rw:coordinates:horizontal-B-margin}
 \partial_{\bar B}H_c
 =1+\bar c\partial_{\bar B}\chi_c-\epsilon\chi_c^2
       +2b\chi_c\partial_{\bar B}\chi_c
 =1+O(u+\epsilon)>1/2.
\end{equation}
Equations \eqref{rw:coordinates:gamma-center} and
\eqref{rw:coordinates:horizontal-B-margin}, together with the common
inverse box and coefficient Lipschitz bounds, will select one weak
comparison center for the whole input interval, independently of
closing zeros.
\end{proof}
